\chapter{Commitments and Secret Sharing}

So far almost every scheme in this book has been about one of two goals: hiding a message from an eavesdropper, or proving who sent it. This chapter is about two quieter primitives that turn out to be load-bearing for an astonishing range of protocols --- electronic voting, auctions, coin-flipping over the phone, threshold wallets, the inner machinery of zero-knowledge proofs and multiparty computation you will meet in the next chapters.

The first is a \emph{commitment}: a way for Alice to lock a value into a sealed envelope, hand the envelope to Bob, and then later open it --- with the guarantees that Bob cannot peek inside before she opens it, and that Alice cannot swap the contents after the fact. The second is \emph{secret sharing}: a way to split a secret into $n$ pieces so that any $t$ of them reconstruct it but any $t-1$ reveal \emph{nothing at all}. Both are built from mathematics you already have --- hash functions, the discrete log problem, and polynomials over a finite field --- and both reward you with that satisfying click of a small idea doing a large job.

\chapterline

\section{The sealed envelope: commitment schemes}
\label{ch32:intro}

Picture a coin flip conducted over the telephone. Alice and Bob are getting a divorce and want to decide who keeps the car. Alice says ``call it'' and Bob says ``heads.'' Alice then announces ``sorry, it was tails.'' Bob has no way to know whether Alice flipped fairly or just declared whatever beat his call. The protocol is hopeless because Alice's commitment to the outcome and her revelation of it happen at the same instant, and she controls both.

A commitment scheme fixes exactly this. Alice flips the coin privately, gets (say) heads, and sends Bob a \emph{commitment} $c$ --- a scrambled token that is bound to ``heads'' but reveals nothing about it. Bob then makes his guess. Only afterward does Alice \emph{open} the commitment, sending Bob the value and whatever extra information he needs to check that $c$ really was a commitment to heads. Now neither party can cheat: Alice cannot change her flip after seeing Bob's guess, and Bob learns nothing from $c$ that would let him guess better than chance. Figure~\ref{ch32:fig-flow} shows the three-message shape.

\begin{intuitionbox}
Think of a commitment as a locked, opaque strongbox. Alice puts her value inside, locks it, and hands the box to Bob (this is \emph{commit}). Bob holds the box but cannot see in (\emph{hiding}). Later Alice hands him the key, he opens it, and reads the value (\emph{open}). Because the box was sealed the whole time, Alice could not have reached in to swap the contents (\emph{binding}). Two properties, two directions of cheating prevented.
\end{intuitionbox}

\begin{figure}[H]
\centering
\begin{tikzpicture}[font=\small]
  \coordinate (A) at (0,0);
  \coordinate (B) at (8,0);
  \node[above,aliceblue] at (A) {\textbf{Alice} (committer)};
  \node[above,bobgreen] at (B) {\textbf{Bob} (receiver)};
  \draw[aliceblue,thick] (A) -- ++(0,-5.2);
  \draw[bobgreen,thick] (B) -- ++(0,-5.2);
  % commit
  \draw[-{Stealth},thick] ($(A)+(0,-1)$) -- node[above,scale=0.85]{commit $c=\operatorname{Commit}(m,r)$} ($(B)+(0,-1)$);
  \node[aliceblue,left,scale=0.8] at ($(A)+(0,-0.5)$){picks $m,r$};
  % bob guesses
  \draw[-{Stealth},thick] ($(B)+(0,-2.2)$) -- node[above,scale=0.85]{Bob acts (guess / bid)} ($(A)+(0,-2.2)$);
  \node[bobgreen,right,scale=0.8] at ($(B)+(0,-1.7)$){$c$ leaks nothing (hiding)};
  % open
  \draw[-{Stealth},thick] ($(A)+(0,-3.6)$) -- node[above,scale=0.85]{open $(m,r)$} ($(B)+(0,-3.6)$);
  \node[bobgreen,right,scale=0.8] at ($(B)+(0,-4.3)$){checks $c\stackrel{?}{=}\operatorname{Commit}(m,r)$};
  \node[aliceblue,left,scale=0.8] at ($(A)+(0,-4.3)$){cannot swap $m$ (binding)};
\end{tikzpicture}
\caption{The commit--reveal flow. Alice binds herself to $m$ with $c$ before Bob acts; the value stays hidden until she opens. Hiding protects Bob's move; binding stops Alice from changing $m$ afterward.}
\label{ch32:fig-flow}
\end{figure}

Formally a commitment scheme is a pair of algorithms. The committing party chooses a message $m$ and a fresh random value $r$ (the \emph{opening} or \emph{nonce}), and computes
\[
  c = \operatorname{Commit}(m, r).
\]
She publishes $c$. To open, she reveals $(m, r)$, and anyone recomputes $\operatorname{Commit}(m,r)$ and checks it equals the $c$ they saw. The two security goals are duals of each other.

\begin{definition}[Hiding and binding]
\label{ch32:hidingbinding}
A commitment scheme is \emph{hiding} if the commitment $c$ leaks nothing about $m$: for any two messages $m_0, m_1$, the distributions of $\operatorname{Commit}(m_0, r)$ and $\operatorname{Commit}(m_1, r)$ over random $r$ are indistinguishable. It is \emph{binding} if the committer cannot open one commitment two ways: it is infeasible to find $(m, r) \neq (m', r')$ with $\operatorname{Commit}(m, r) = \operatorname{Commit}(m', r')$.
\end{definition}

Each property can hold \emph{perfectly} (against an unbounded adversary) or \emph{computationally} (against a $\PPT$ adversary). A short but deep fact: you cannot have both perfectly at once.

\begin{whybox}
Why can't a commitment be both perfectly hiding \emph{and} perfectly binding? Suppose it were perfectly hiding. Then for every commitment $c$ and every message $m$, there must exist some opening $r$ with $\operatorname{Commit}(m, r) = c$ --- otherwise $c$ would rule out $m$, leaking information, contradicting perfect hiding. But if every message has a valid opening of $c$, an unbounded Alice can find an opening to \emph{any} message she likes, so the scheme cannot be binding against an unbounded adversary. The contrapositive runs the same way. So one of the two properties is always merely computational. Designers choose which side to make unconditional based on the threat model: a scheme that must protect a secret \emph{forever} wants perfect hiding; one that must prevent forgery against an all-powerful prover wants perfect binding.
\end{whybox}

\subsection{The threat model}
\label{ch32:threatmodel}

Spell out who is trying to do what. Against \emph{hiding}, the adversary is the receiver Bob, who holds $c$ and wants to learn $m$ before Alice opens. Against \emph{binding}, the adversary is the sender Alice, who has committed to $c$ and now wants to open it as a different value than she originally intended. A good scheme defends against both, and the security argument must say which party is bounded and which capability is assumed hard.

\section{Building commitments from hash functions}
\label{ch32:hashcommit}

The most direct construction uses a cryptographic hash. Let $\Hash$ be a hash function (think SHA-256). To commit to a message $m$, pick a long random nonce $r$ --- say 256 bits --- and publish
\[
  c = \Hash(m \concat r).
\]
To open, reveal $(m, r)$; the verifier recomputes $\Hash(m \concat r)$ and checks it matches $c$.

\begin{whybox}
\textbf{Binding} reduces to collision resistance. If Alice could open $c$ as both $(m,r)$ and $(m',r')$ with $m \ne m'$, then $\Hash(m\concat r) = \Hash(m'\concat r')$ is a collision in $\Hash$ --- which we assume is infeasible to find. So binding holds as long as the hash is collision resistant.

\textbf{Hiding} comes from the nonce. Without $r$, committing to a coin flip as $\Hash(\text{``heads''})$ would be a disaster: Bob just computes $\Hash(\text{``heads''})$ and $\Hash(\text{``tails''})$ himself and compares. The message space is tiny, so the commitment is trivially invertible by brute force. The random $r$ blows the effective input space up to $2^{256}$ possibilities, so Bob cannot enumerate. Formally, hiding here is computational and is usually argued by modeling $\Hash$ as a random oracle, in which case $c$ is just a uniformly random string independent of $m$ until someone queries the exact $m\concat r$.
\end{whybox}

\begin{trapbox}
The nonce is not optional, and it must be \emph{long and unpredictable}. Two real failures: (1) committing to a low-entropy value with no nonce (a vote, a bid, a yes/no) lets the receiver brute-force the message in microseconds. (2) Reusing the same $r$ across commitments, or deriving $r$ from $m$ deterministically, can leak relationships between committed messages. Always draw $r$ fresh and uniformly. A commitment to a 1-bit secret with a 256-bit random nonce is fine; a commitment to the same bit with no nonce is worthless.
\end{trapbox}

\begin{implementbox}
A toy hash commitment in Python. This is for \emph{learning only} --- it skips constant-time comparison, input-length framing (so $m\concat r$ is unambiguous), and the domain separation a real library would insist on.
\begin{lstlisting}[language=Python]
import hashlib, os

def commit(message: bytes):
    r = os.urandom(32)                      # 256-bit nonce, fresh every time
    c = hashlib.sha256(message + r).digest()
    return c, r                             # publish c, keep (message, r) to open

def verify(c, message: bytes, r: bytes) -> bool:
    return hashlib.sha256(message + r).digest() == c

# Alice flips a coin and commits.
flip = b"heads"
c, r = commit(flip)
print("commitment:", c.hex())     # Bob sees only this; learns nothing about flip

# ... Bob makes his guess ...

# Alice opens.
print("opens to:", flip, "valid?", verify(c, flip, r))   # True
# Alice cannot now claim it was b"tails": verify(c, b"tails", r) is False,
# and finding r' with sha256(b"tails"+r')==c is a SHA-256 preimage break.
\end{lstlisting}
A subtle bug lurks in \texttt{message + r}: if message lengths can vary, \texttt{b"ab"+b"c"} and \texttt{b"a"+b"bc"} collide. Real code length-prefixes each field. We will let it slide for a fixed-length flip.
\end{implementbox}

The hash commitment is computationally hiding and computationally binding. That is enough for most applications, and it is very fast. But sometimes you want stronger guarantees or extra structure --- in particular, the ability to do arithmetic \emph{on commitments without opening them}. For that we turn to the discrete log.

\section{Pedersen commitments: perfect hiding and homomorphism}
\label{ch32:pedersen}

Let $\G$ be a cyclic group of prime order $q$ where the discrete log problem is hard --- a subgroup of $\Zns{p}$, or an elliptic curve group from Chapter 21. Fix two generators $g$ and $h$ whose relative discrete log nobody knows: $h = g^{x}$ for some $x$ that not even the person who set up the parameters has stored. (One produces $h$ by hashing a fixed string to a group element, so that $x$ is unknown to everyone.) To commit to a message $m \in \Zn{q}$, pick a random $r \in \Zn{q}$ and publish
\[
  c = g^{m} h^{r}.
\]
To open, reveal $(m, r)$; the verifier checks $c \stackrel{?}{=} g^m h^r$.

\begin{whybox}
\textbf{Perfectly hiding.} For any fixed $m$, as $r$ ranges uniformly over $\Zn{q}$, the term $h^r$ ranges uniformly over all of $\G$ (because $h$ generates the group and $q$ is prime). So $c = g^m h^r$ is a uniformly random group element, with a distribution that does not depend on $m$ in the slightest. Given $c$, \emph{every} message $m'$ is equally consistent: there is exactly one $r'$ making $c = g^{m'} h^{r'}$. An unbounded Bob learns nothing. This is information-theoretic hiding --- it holds forever, even against a quantum computer.

\textbf{Computationally binding.} Suppose Alice opens $c$ two ways: $g^{m}h^{r} = g^{m'}h^{r'}$ with $m \ne m'$. Rearranging, $g^{\,m-m'} = h^{\,r'-r}$, so
\[
  h = g^{(m-m')(r'-r)^{-1} \bmod q},
\]
and Alice has just computed $\dlog_g h = x$, the one number nobody is supposed to know. So breaking binding is exactly as hard as solving discrete log for $h$ relative to $g$. As long as that is infeasible, Alice is stuck with her commitment.
\end{whybox}

Compare the two schemes. The hash version is computationally hiding and computationally binding. The Pedersen version trades up to \emph{perfect} hiding at the cost of making binding rest on discrete log (and being far slower, since it needs group exponentiations). By the impossibility result, you cannot get perfect on both sides, and Pedersen has chosen the hiding side. If you instead want perfect binding, use $c = g^m$ with no nonce in a group where $g$ has prime order --- now each $m$ gives a unique $c$, but the scheme is only computationally hiding and only safe for high-entropy $m$.

\subsection{The homomorphic bonus}
\label{ch32:homomorphic}

Pedersen commitments have a property the hash version cannot match: they add. Multiply two commitments and the exponents add,
\[
  c_1 \cdot c_2 = (g^{m_1} h^{r_1})(g^{m_2} h^{r_2}) = g^{\,m_1+m_2}\, h^{\,r_1+r_2},
\]
which is a valid commitment to $m_1 + m_2$ with opening $r_1 + r_2$. So you can prove statements about \emph{sums} of committed values without ever opening the individual values. A confidential transaction can commit to each amount, and anyone can verify that inputs equal outputs --- that no money was created --- by checking that the product of input commitments equals the product of output commitments, all while the amounts themselves stay hidden. This is the backbone of confidential transactions in cryptocurrencies and a building block for the multiparty computation of Chapter 34.

\begin{walkbox}
\emph{Question.} Suppose Alice commits to $m_1 = 7$ and $m_2 = 11$ with the same nonce $r$ in both, instead of fresh nonces. What goes wrong?

\emph{Answer.} Hiding within each commitment still holds, but the \emph{relationship} leaks through the homomorphism. Consider $c_1 / c_2 = g^{m_1 - m_2} h^{r - r} = g^{m_1 - m_2} = g^{-4}$. The nonces cancel, leaving a nonce-free commitment $g^{m_1-m_2}$ that is only computationally hiding --- and if the difference comes from a small set, Bob brute-forces it. The lesson is the same as for hash commitments: nonces must be fresh and independent.
\end{walkbox}

\begin{applicationbox}
Commitments are everywhere once you look. \textbf{TLS and Signal} use the hash-commit pattern implicitly whenever a protocol transcript is hashed and bound before later steps. \textbf{Certificate Transparency} logs commit to the set of issued certificates via Merkle trees (a tree of hash commitments) so that a misissued certificate cannot be hidden. \textbf{Cryptocurrencies} use Pedersen commitments for confidential amounts (Monero, Mimblewimble) and hash commitments in commit-reveal schemes for on-chain randomness and sealed-bid auctions. \textbf{Zero-knowledge proofs} (Chapter 31) are riddled with commitments: the prover commits to its randomness, the verifier challenges, the prover opens --- the commit-challenge-open shape is the skeleton of nearly every interactive proof.
\end{applicationbox}

\section{Splitting a secret: the problem}
\label{ch32:ss-problem}

Now a different scenario. The launch code for a missile silo, or the master key that decrypts a company's password vault, is too dangerous to entrust to any single person and too important to lose if one person is hit by a bus. You want to hand pieces to five officers such that any \emph{three} of them together can reconstruct the key, but any two --- even two who collude, even two with unlimited computing power --- learn \emph{absolutely nothing}. This is a $(t,n)$ \emph{threshold secret sharing} scheme with $t=3$ and $n=5$.

The naive ideas all fail. Giving each officer a third of the bits of the key? Then two officers know two-thirds of the bits and can brute-force the rest --- and any officer alone already knows a third, which is a leak. Encrypting the key and splitting the ciphertext? Same problem; partial ciphertext is partial information. We want a scheme where $t-1$ shares are \emph{statistically independent} of the secret, so they reveal nothing in the information-theoretic sense, not merely the computational one.

\begin{intuitionbox}
The trick rests on a fact from high-school geometry, sharpened. Two points determine a unique line. If I draw a line and the secret is its height where it crosses the $y$-axis, then any \emph{one} point on the line tells you nothing about that height --- infinitely many lines pass through a single point, each hitting the axis somewhere different. But \emph{two} points pin the line down completely, and you can read off the intercept. To go from ``2 reconstruct'' to ``$t$ reconstruct,'' replace the line (degree 1) with a polynomial of degree $t-1$: it takes exactly $t$ points to determine it, and any $t-1$ points leave it completely undetermined.
\end{intuitionbox}

\section{Shamir secret sharing}
\label{ch32:shamir}

Adi Shamir (the S in RSA) published this in 1979 in a two-page paper titled ``How to Share a Secret.'' It is one of the most elegant constructions in all of cryptography, and it is nothing but polynomial interpolation over a finite field.

\begin{historybox}
Shamir's paper and George Blakley's geometric scheme appeared independently the same year (1979). Blakley used intersecting hyperplanes; Shamir used polynomials. Shamir's is the one that stuck, because shares are small (one field element each, regardless of $t$) and the math is clean. The same year gave us a remarkable cluster of foundational ideas --- it sits right after the 1976--1978 burst of Diffie-Hellman and RSA, as the young field was working out what public mathematics could do for secrets.
\end{historybox}

\subsection{Setup and why a field}
\label{ch32:field}

All arithmetic happens in a finite field $\F_p$ for a prime $p$ larger than both the secret $s$ and the number of participants $n$. (One can also use $\F_{2^k}$ for byte-aligned implementations; the algebra is identical.) We need a field --- not just integers mod a composite --- because reconstruction divides by differences of evaluation points, and division is only guaranteed in a field, where every nonzero element has an inverse.

\begin{connectionbox}
This is the payoff of the finite-field chapter (Chapter 6). A field is precisely the structure where the four operations $+,-,\times,\div$ all behave, and Lagrange interpolation needs all four. Over $\Zn{n}$ with composite $n$ you can hit a non-invertible difference $x_i - x_j$ and the whole reconstruction collapses. Over $\F_p$ it never happens, as long as the evaluation points are distinct and nonzero mod $p$. The companion algorithms book's treatment of polynomial evaluation and the cost of modular inversion (via extended Euclid) is exactly the cost model for sharing and reconstruction here.
\end{connectionbox}

\subsection{Dealing the shares}
\label{ch32:deal}

To share a secret $s \in \F_p$ among $n$ players with threshold $t$:

\begin{enumerate}
  \item Set the constant term to the secret: $a_0 = s$.
  \item Pick the other coefficients $a_1, \dots, a_{t-1}$ \emph{uniformly at random} from $\F_p$. This defines a random polynomial of degree $t-1$:
  \[
    f(x) = a_0 + a_1 x + a_2 x^2 + \dots + a_{t-1}x^{t-1} \pmod p .
  \]
  \item Give player $i$ (for $i = 1, \dots, n$) the share $(i, f(i))$. The $x$-coordinates are public labels; the $y$-values are the secret shares. Crucially, \emph{nobody gets} $f(0) = s$ itself --- that point is the secret.
\end{enumerate}

To reconstruct from any $t$ shares, interpolate the unique degree-$(t-1)$ polynomial through them and evaluate at $0$.

\subsection{A worked example over $\F_{17}$}
\label{ch32:worked}

Let me share the secret $s = 5$ with threshold $t = 3$ among $n = 5$ players, working in $\F_{17}$ (so $p = 17 > n$ and $> s$). Because $t = 3$ I need a degree-2 polynomial. Set $a_0 = 5$ and pick random $a_1 = 8$, $a_2 = 3$:
\[
  f(x) = 5 + 8x + 3x^2 \pmod{17}.
\]
Now evaluate at $x = 1, \dots, 5$, reducing mod 17 each time:
\[
\begin{aligned}
f(1) &= 5 + 8 + 3 = 16, \\
f(2) &= 5 + 16 + 12 = 33 \equiv 16, \\
f(3) &= 5 + 24 + 27 = 56 \equiv 56 - 51 = 5, \\
f(4) &= 5 + 32 + 48 = 85 \equiv 85 - 85 = 0, \\
f(5) &= 5 + 40 + 75 = 120 \equiv 120 - 119 = 1.
\end{aligned}
\]
So the five shares are
\[
  (1,16),\ (2,16),\ (3,5),\ (4,0),\ (5,1).
\]
Hand one to each officer. Note $(4,0)$: a share can perfectly well be zero, and that zero leaks nothing.

\begin{center}
\begin{tikzpicture}[scale=0.42]
  \draw[->,gray] (-0.5,0) -- (6.5,0) node[right,black]{$x$};
  \draw[->,gray] (0,-0.5) -- (0,17.5) node[above,black]{$f(x)$};
  \foreach \x in {1,2,3,4,5} \draw (\x,0.15)--(\x,-0.15) node[below,scale=0.7]{\x};
  \foreach \y in {10,15} \draw (0.15,\y)--(-0.15,\y) node[left,scale=0.7]{\y};
  % secret point at x=0
  \fill[evered] (0,5) circle (6pt) node[left=2pt,black,scale=0.8]{$s=5$};
  % share points
  \fill[aliceblue] (1,16) circle (5pt);
  \fill[aliceblue] (2,16) circle (5pt);
  \fill[aliceblue] (3,5) circle (5pt);
  \fill[aliceblue] (4,0) circle (5pt);
  \fill[aliceblue] (5,1) circle (5pt);
  \node[scale=0.7,above right] at (1,16) {(1,16)};
  \node[scale=0.7,above right] at (5,1) {(5,1)};
\end{tikzpicture}
\end{center}

The picture is over the integers for intuition only --- the real points live wrapped around $\F_{17}$ --- but it conveys the idea: the secret is the hidden intercept $f(0)$, and the shares are scattered points on a curve only $t$ of them can identify.

\subsection{Reconstruction by Lagrange interpolation}
\label{ch32:reconstruct}

Take any three shares --- say officers 2, 3, and 5 hold $(2,16), (3,5), (5,1)$. The unique degree-2 polynomial through them is given by the Lagrange formula. We do not even need the whole polynomial; we only want $f(0) = s$. The Lagrange basis evaluated at $x=0$ gives
\[
  s = f(0) = \sum_{j} y_j \prod_{m \ne j} \frac{0 - x_m}{x_j - x_m} \pmod{p},
\]
where the sum runs over the chosen shares. Compute each coefficient $\lambda_j = \prod_{m\ne j} \frac{-x_m}{x_j - x_m}$ for our points $x = 2,3,5$.

For $j$ at $x_j = 2$:
\[
  \lambda_2 = \frac{(0-3)(0-5)}{(2-3)(2-5)} = \frac{(-3)(-5)}{(-1)(-3)} = \frac{15}{3} = 5.
\]
For $j$ at $x_j = 3$:
\[
  \lambda_3 = \frac{(0-2)(0-5)}{(3-2)(3-5)} = \frac{(-2)(-5)}{(1)(-2)} = \frac{10}{-2} = -5 \equiv 12 \pmod{17}.
\]
For $j$ at $x_j = 5$:
\[
  \lambda_5 = \frac{(0-2)(0-3)}{(5-2)(5-3)} = \frac{(-2)(-3)}{(3)(2)} = \frac{6}{6} = 1.
\]
Now combine with the $y$-values $16, 5, 1$:
\[
  s = 16\cdot 5 + 5\cdot 12 + 1\cdot 1 = 80 + 60 + 1 = 141 \equiv 141 - 8\cdot 17 = 141 - 136 = 5 \pmod{17}.
\]
We recover $s = 5$. And it works no matter \emph{which} three shares we pick --- try officers 1, 4, 5 and you will land on 5 again.

\begin{walkbox}
\emph{Question.} In computing $\lambda_3$ we wrote $\frac{10}{-2}$. Over the rationals that is $-5$; over $\F_{17}$ what does division even mean here, and why did $-5$ become $12$?

\emph{Answer.} Division by $-2$ in $\F_{17}$ means multiplication by the inverse of $-2 \equiv 15$. Since $15 \cdot 8 = 120 = 7\cdot 17 + 1 \equiv 1$, we have $15^{-1} = 8$, so $10 \cdot 8 = 80 \equiv 80 - 68 = 12$. And indeed $-5 \equiv 12 \pmod{17}$. The two routes agree, which is the whole point of working in a field: the rational arithmetic and the modular arithmetic are the same operation. This is also where a composite modulus would bite --- if $x_j - x_m$ shared a factor with the modulus, no inverse would exist and the formula would jam.
\end{walkbox}

\subsection{Why $t-1$ shares reveal nothing}
\label{ch32:perfect}

This is the heart of the scheme, and it is genuinely perfect secrecy, not a computational assumption.

\begin{theorem}[Information-theoretic security of Shamir]
\label{ch32:thm-perfect}
Given any $t-1$ shares, every candidate secret $s' \in \F_p$ is equally likely. The shares are statistically independent of the secret.
\end{theorem}

\begin{whybox}
Fix any $t-1$ shares $(x_1,y_1),\dots,(x_{t-1},y_{t-1})$ and propose any secret value $s'$. Claim: there is \emph{exactly one} degree-$(t-1)$ polynomial consistent with both the $t-1$ shares and the additional point $(0, s')$. That is $t$ points total, and $t$ points determine a unique degree-$(t-1)$ polynomial --- so the candidate $(0,s')$ together with the held shares pins down exactly one polynomial, for \emph{every} choice of $s'$.

Because the dealer chose the coefficients $a_1,\dots,a_{t-1}$ uniformly at random, each of those $p$ possible polynomials (one per value of $s'$) was equally likely to have been the one used. So conditioned on the $t-1$ observed shares, the posterior distribution over the secret is uniform on $\F_p$: the shares told the attacker nothing. This is the same flavor of argument as the one-time pad in Chapter 9 --- every plaintext stays equally consistent with what the adversary sees.
\end{whybox}

Concretely with our example: hold the two shares $(1,16),(2,16)$. Could the secret be $5$? Yes --- our actual polynomial works. Could it be $0$? Also yes: there is a unique degree-2 polynomial through $(0,0),(1,16),(2,16)$, namely some other $f$ with $a_0 = 0$, equally plausible. Two shares cannot tell these apart. Only the third point breaks the tie.

\begin{implementbox}
Shamir sharing and reconstruction in Python over a prime field. \emph{Learning only}: a real library uses a constant-time field, validates share indices, authenticates shares (see verifiable sharing below), and uses a cryptographically secure RNG with care. The standard library \texttt{random} below is \textbf{not} secure --- swap in \texttt{secrets} for anything real.
\begin{lstlisting}[language=Python]
import random
P = 2**127 - 1   # a Mersenne prime field; secret must be < P

def share(secret, t, n):
    # coefficients: a0 = secret, rest uniform in F_P
    coeffs = [secret] + [random.randrange(P) for _ in range(t - 1)]
    def f(x):
        y = 0
        for a in reversed(coeffs):     # Horner's rule
            y = (y * x + a) % P
        return y
    return [(i, f(i)) for i in range(1, n + 1)]

def reconstruct(shares):
    # Lagrange interpolation evaluated at x = 0
    secret = 0
    for j, (xj, yj) in enumerate(shares):
        num, den = 1, 1
        for m, (xm, _) in enumerate(shares):
            if m == j:
                continue
            num = (num * (-xm)) % P          # (0 - xm)
            den = (den * (xj - xm)) % P      # (xj - xm)
        lam = num * pow(den, -1, P) % P      # modular inverse via Fermat/EEA
        secret = (secret + yj * lam) % P
    return secret

shares = share(1234567, t=3, n=5)
print(reconstruct(shares[:3]))    # 1234567, any 3 of the 5 work
print(reconstruct(shares[1:4]))   # 1234567 again
\end{lstlisting}
\end{implementbox}

\subsection{Properties that make it shine}
\label{ch32:properties}

A few features fall out for free and explain why Shamir is the default in practice:

\begin{itemize}
  \item \textbf{Minimal share size.} Each share is one field element, the same size as the secret, regardless of $t$. You pay no penalty for a high threshold.
  \item \textbf{Flexible threshold.} Choose any $(t,n)$ with $t \le n$. New players can be added by handing out $f(n+1), f(n+2), \dots$ without disturbing existing shares.
  \item \textbf{Proactive refresh.} The dealer (or the players jointly) can add a fresh random polynomial $g$ with $g(0) = 0$ to everyone's share: $f(i) \mapsto f(i) + g(i)$. The secret $f(0)+g(0) = f(0)$ is unchanged, but old shares are now useless. This periodically ``re-randomizes'' shares so an attacker must compromise $t$ players within a single epoch.
  \item \textbf{Weighted access.} Give an important player several shares (several $x$-coordinates) to give them more voting weight. A vice-president with two shares plus one officer can hit $t=3$.
\end{itemize}

\section{The dealer problem and verifiable secret sharing}
\label{ch32:vss}

Plain Shamir assumes a trustworthy dealer and cooperative players. Two ways that breaks down in the wild:

\begin{attackbox}
\textbf{A cheating dealer.} A malicious dealer can hand out inconsistent shares --- points that do not lie on any single degree-$(t-1)$ polynomial. Then different subsets of $t$ players reconstruct \emph{different} secrets, and nobody can tell who got bad data. In our $\F_{17}$ example, suppose the dealer gives player 3 the share $(3, 9)$ instead of the correct $(3,5)$. Then officers $\{2,3,5\}$ interpolate one value and officers $\{1,2,4\}$ interpolate another. The shares are individually plausible, so the fraud is invisible until reconstruction --- by which time you may have lost the only chance to fix it.

\textbf{A cheating player at reconstruction.} Even with a trustworthy dealer, a single player can submit a forged share during reconstruction and steer the recovered value to anything they like (since they control one Lagrange term). Worse, they can submit the forgery, watch the others reveal genuine shares, and then privately recover the true secret while everyone else gets garbage.
\end{attackbox}

\emph{Verifiable secret sharing} (VSS) closes both holes by publishing commitments that let every player check their own share against a public, binding record of the polynomial --- without revealing the secret. Feldman's 1987 scheme is the classic, and it reuses the discrete-log commitment machinery from earlier in this chapter.

\subsection{Feldman VSS}
\label{ch32:feldman}

Work in a group $\G$ of prime order $q$ with generator $g$ where discrete log is hard, and take the secret and coefficients in $\F_q$. The dealer runs ordinary Shamir to build $f(x) = a_0 + a_1 x + \dots + a_{t-1}x^{t-1}$ with $a_0 = s$, and distributes shares $f(i)$ privately as usual. \emph{In addition}, the dealer broadcasts public commitments to every coefficient:
\[
  C_0 = g^{a_0},\quad C_1 = g^{a_1},\quad \dots,\quad C_{t-1} = g^{a_{t-1}}.
\]
Now any player $i$ can verify their share $f(i)$ is consistent with the public commitments, because exponentiation turns the polynomial identity into a checkable group equation:
\[
  g^{f(i)} \stackrel{?}{=} \prod_{k=0}^{t-1} C_k^{\,i^{k}} = g^{a_0}\,(g^{a_1})^{i}\,(g^{a_2})^{i^2}\cdots .
\]

\begin{whybox}
Why does the product equal $g^{f(i)}$? Just collect exponents:
\[
  \prod_{k=0}^{t-1} C_k^{\,i^k} = \prod_k (g^{a_k})^{i^k} = g^{\sum_k a_k i^k} = g^{f(i)}.
\]
So a share passes the check exactly when it sits on the polynomial whose coefficients the dealer committed to. The dealer cannot hand out an off-curve share to anyone without it failing this public test --- binding of the commitments (discrete log hardness) forbids opening $C_k$ to a different $a_k$. And the commitments leak nothing usable about $s$ when $s$ is a high-entropy key: $C_0 = g^{s}$ is just a discrete-log-hard image of the secret, computationally hiding.
\end{whybox}

\begin{trapbox}
Feldman VSS hides the secret only \emph{computationally}, because $C_0 = g^s$ commits to $s$ with no nonce. If the secret has low entropy, $g^s$ can be brute-forced --- exactly the no-nonce pitfall from the hash commitment. Pedersen VSS fixes this by committing to each coefficient with a paired blinding polynomial, $C_k = g^{a_k} h^{b_k}$, recovering perfect hiding of the secret at the cost of one extra polynomial. Use Feldman when the secret is itself a high-entropy key; use Pedersen VSS when the secret might be guessable.
\end{trapbox}

\section{Distributed key generation: removing the dealer}
\label{ch32:dkg}

VSS still has \emph{one} dealer who knows the whole secret at deal time. For the highest-assurance systems --- a threshold signing key that should never exist in one place, ever --- you want \emph{distributed key generation} (DKG). The idea composes the pieces we have:

\begin{enumerate}
  \item Every one of the $n$ players acts as a dealer for a random secret $s_i$ of their own, running Feldman VSS to share $s_i$ to everyone and publishing coefficient commitments.
  \item Each player verifies the shares they received from every other player; complaints disqualify cheaters.
  \item The group secret is defined as $s = \sum_i s_i$ over the surviving dealers, and each player's share of $s$ is the sum of the shares they received, $\sum_i f_i(\text{me})$.
\end{enumerate}

By the homomorphism of Shamir sharing (the sum of polynomials sharing $s_i$ is a polynomial sharing $\sum s_i$), the players now hold valid shares of a secret that \emph{no single participant ever knew} --- not even the dealers, since $s$ depends on everyone's contribution. The public key $g^s = \prod_i C_0^{(i)}$ is computable by all, while the private key $s$ lives only as a threshold sharing.

\begin{applicationbox}
DKG plus threshold signatures is how modern \textbf{threshold wallets} and \textbf{validator key management} work. A custody service can require $t$-of-$n$ hardware modules to cooperate to sign a transaction, so stealing fewer than $t$ devices yields nothing and losing up to $n-t$ devices loses no funds. The same machinery secures \textbf{distributed certificate authorities}, \textbf{DNSSEC root key ceremonies} (the literal ceremony in a guarded room with smart cards is a physical instantiation of $(t,n)$ sharing), and \textbf{validator clusters} in proof-of-stake blockchains that sign as a group. Threshold ECDSA and threshold BLS signatures, both built on these primitives, run in production custody systems today.
\end{applicationbox}

\section{Where this gets used: key escrow, multisig, distributed trust}
\label{ch32:applications}

Step back and look at the family of problems secret sharing solves. They all share a shape: \emph{no single point of failure, no single point of trust.}

\textbf{Key escrow and recovery.} An enterprise encrypts its data with a master key. If the key is lost the data is gone; if one administrator holds it, that administrator is a liability and a target. Split the master key $3$-of-$5$ across executives kept in separate locations. Day-to-day operations use a separate operational key; the master is reconstructed only in a genuine recovery event, and only by a quorum. The infamous government ``key escrow'' proposals of the 1990s (the Clipper chip) used related ideas to give law enforcement a recovery path --- and the fierce debate over them is a reminder that the policy question (who should be in the quorum?) is harder than the math.

\textbf{Multisignature wallets.} ``Multisig'' in cryptocurrencies often means $t$-of-$n$ \emph{independent} signatures recorded on chain (a script checks that $t$ distinct public keys signed). Threshold signatures via DKG are the cryptographic cousin: they produce a \emph{single} signature under a single public key that is indistinguishable from an ordinary one, but which required $t$ parties to cooperate. The on-chain version is transparent and simple; the threshold version is private and compact. Both deliver the same governance: a $2$-of-$3$ wallet survives losing one key and resists one compromised key.

\textbf{Distributed trust.} Certificate authorities, code-signing roots, and password vaults all concentrate enormous trust in one secret. Sharing it across mutually-suspicious parties, in separate jurisdictions or organizations, means an attacker must subvert a quorum rather than a single box --- and proactive refresh means they must do it fast. This is defense in depth applied to the secret itself.

\begin{connectionbox}
Secret sharing is the gateway to the two chapters ahead. \textbf{Multiparty computation} (Chapter 34) lets the shareholders compute on the shared secret --- adding, even multiplying, shared values --- so that a function of private inputs is evaluated without anyone seeing the inputs. The homomorphism we used in DKG ($\sum$ of shares = share of $\sum$) is the addition gate of MPC; multiplication needs more work but starts here. And \textbf{zero-knowledge proofs} (Chapter 31) lean on commitments at every turn. The two primitives in this chapter are not side curiosities; they are the connective tissue of advanced cryptographic protocols.
\end{connectionbox}

\section{Summary}
\label{ch32:summary}

A commitment is a cryptographic sealed envelope: \emph{hiding} keeps the receiver from peeking, \emph{binding} keeps the sender from swapping. You cannot have both unconditionally, so you pick a side. Hash commitments $\Hash(m\concat r)$ are fast and computationally secure on both fronts --- never forget the nonce $r$ for low-entropy messages. Pedersen commitments $g^m h^r$ buy perfect hiding from the discrete log assumption and throw in an additive homomorphism that powers confidential transactions.

Shamir secret sharing splits a secret into $n$ shares so any $t$ reconstruct it and any $t-1$ learn nothing --- perfect, information-theoretic secrecy --- using a random degree-$(t-1)$ polynomial whose constant term is the secret. Reconstruction is Lagrange interpolation in a finite field. Verifiable secret sharing (Feldman, Pedersen) adds public commitments so players can detect a cheating dealer; distributed key generation removes the dealer entirely, producing a key no one ever held. These power key escrow, threshold wallets, and distributed trust --- and they are the foundation for the MPC and zero-knowledge protocols still to come.

\section{Exercises}
\label{ch32:exercises}

\begin{exercise}
\emph{(Compute a hash commitment by hand-ish.)} Alice commits to the bit $b = 1$ with no nonce, publishing $c = \Hash(\text{``1''})$. Bob wants to learn $b$ before Alice opens. Describe exactly what Bob computes and how many hash evaluations it costs him. Then redo the analysis when Alice uses a uniform 256-bit nonce, and explain why Bob's strategy now fails.
\end{exercise}

\begin{exercise}
\emph{(Binding break is a collision.)} Suppose someone hands you an efficient algorithm that, given a hash commitment $c = \Hash(m\concat r)$, outputs $(m',r') \ne (m,r)$ that also opens $c$. Show step by step how to turn this into a collision finder for $\Hash$. Conclude that a collision-resistant hash yields a binding commitment.
\end{exercise}

\begin{exercise}
\emph{(Pedersen perfect hiding.)} In a group of prime order $q$ with generators $g,h$, show that for a fixed message $m$ the commitment $c = g^m h^r$ is uniformly distributed over the whole group as $r$ ranges uniformly over $\F_q$. Then show that for any \emph{other} message $m'$ there is exactly one $r'$ with $g^{m'}h^{r'} = c$, and write $r'$ in terms of $m, m', r$, and $x = \dlog_g h$.
\end{exercise}

\begin{exercise}
\emph{(Equivocation reveals the trapdoor.)} Prove that if Alice can open a single Pedersen commitment $c$ as two different messages, she can compute $\dlog_g h$. Give the explicit formula. Why does this mean binding is exactly as hard as discrete log?
\end{exercise}

\begin{exercise}
\emph{(Share a secret yourself.)} Working in $\F_{23}$, share the secret $s = 9$ with threshold $t = 2$ and $n = 4$ using the line $f(x) = 9 + 7x \bmod 23$. List the four shares. Then take any two of them and reconstruct $s$ via Lagrange interpolation, showing the modular inverse you use.
\end{exercise}

\begin{exercise}
\emph{(Threshold three.)} In $\F_{17}$, the polynomial $f(x) = 5 + 8x + 3x^2$ from the chapter produced shares $(1,16),(2,16),(3,5),(4,0),(5,1)$. Reconstruct $s$ from the shares of officers $\{1,4,5\}$ and confirm you get $5$. Then show that the two shares $\{1,4\}$ alone are consistent with the secret being any value in $\F_{17}$ by exhibiting, for the candidate $s' = 0$, the unique degree-2 polynomial through $(0,0),(1,16),(4,0)$.
\end{exercise}

\begin{exercise}
\emph{(Why a field, not a ring.)} Suppose you tried to run Shamir over $\Zn{12}$ (composite) with evaluation points $x = 1,\dots,n$. Find two evaluation points whose difference is not invertible mod $12$, and explain concretely how Lagrange reconstruction can fail. What property of $\F_p$ guarantees this never happens?
\end{exercise}

\begin{exercise}
\emph{(Proactive refresh.)} Players hold Shamir shares of a secret $s$ under a degree-2 polynomial. The group jointly generates a random degree-2 polynomial $g$ with $g(0) = 0$ and each player adds $g(i)$ to their share. Prove the secret is unchanged, that old and new shares cannot be mixed to reconstruct, and explain in one sentence why this defends against a slow attacker compromising players over time.
\end{exercise}

\begin{exercise}
\emph{(Catch the cheating dealer.)} In Feldman VSS over a prime-order group, the dealer publishes $C_0,\dots,C_{t-1}$ and gives player $i$ a share $y_i$. Write the exact equation player $i$ checks. Then suppose a dealer sends player 3 a tampered share $y_3' = y_3 + 1$; show that the verification equation fails, and argue why the dealer cannot adjust the public commitments to make it pass without breaking discrete log.
\end{exercise}

\begin{exercise}
\emph{(Homomorphic shares.)} Let $f$ share secret $s$ and $g$ share secret $s'$, both degree $t-1$ with the same evaluation points. Show that the pointwise sums $f(i)+g(i)$ are valid shares of $s + s'$ under a degree-$(t-1)$ polynomial. Explain how distributed key generation uses exactly this fact so that no single dealer knows the final key.
\end{exercise}

\begin{exercise}
\emph{(Weighted threshold design.)} You must build a policy where the CEO alone plus any one of four VPs can recover a key, OR any three VPs together can recover it (without the CEO). Using Shamir with a suitable threshold and by handing some parties multiple share-indices, design a concrete scheme: state the threshold $t$ and how many share-indices each role holds, and verify both authorized patterns reach $t$ while a single VP alone does not.
\end{exercise}
